% === III. METHODOLOGY ===
\section{Methodology}
\label{sec:methodology}

\subsection{System Architecture}
\label{subsec:architecture}

\system{} follows a five-stage pipeline architecture (Fig.~\ref{fig:architecture}). The design prioritizes modularity, whole-body coverage, and real-time performance on consumer hardware.

\begin{enumerate}[itemsep=2pt, topsep=3pt, parsep=0pt]
    \item \textbf{Sensing \& Inputs}: Webcam capture (30~FPS), reference exercise template (angular dataset), and Safe-Max ROM calibration from user profile (age, ROM limits, medical conditions).
    \item \textbf{Perception}: Whole-body 3D pose estimation via RTMW3D-L~\cite{ref_rtmw3d} (133 keypoints: body + hands + face), face analysis via OpenFace~3.0 (8 AUs + gaze), and Whisper ASR for speech input.
    \item \textbf{Kinematics \& Alignment}: Quaternion-based joint angle computation, 4th-order Butterworth filtering (6.0~Hz cutoff), and Sakoe-Chiba constrained DTW alignment.
    \item \textbf{Evaluation}: Six-dimension scoring (ROM, stability, flow, symmetry, compensation, smoothness), compensation detection (shoulder hiking, trunk lean, hip shift), and multi-indicator fatigue analysis.
    \item \textbf{Coaching \& HUD}: LLM coaching with RAG grounding, safety guardrails, Edge-TTS voice synthesis, and visual overlay (skeleton with angle arcs, real-time metrics dashboard, session reports).
\end{enumerate}

\begin{figure*}[t]
    \centering
    % =============================================================================
%  ADAPT-Rehab -- Figure 1: Horizontal System Workflow Diagram (body)
%  Input via % =============================================================================
%  ADAPT-Rehab -- Figure 1: Horizontal System Workflow Diagram (body)
%  Input via % =============================================================================
%  ADAPT-Rehab -- Figure 1: Horizontal System Workflow Diagram (body)
%  Input via \input{figures/fig1_architecture_body.tex} into IEEEtran figure*
%  Total visual width ≈ 18.1 cm with scale=0.92 (fits IEEE 2-col textwidth)
% =============================================================================

% --- Dimension constants (cm) ------------------------------------------------
\def\colW{3.5}
\def\colH{7.8}
\def\colGap{0.55}
\def\boxW{3.1}
\def\boxH{1.0}
\def\titleH{0.6}

% --- Component Box Helper Macro ----------------------------------------------
\newcommand{\compbox}[5][]{%
  \node[rounded corners=4pt, draw=NeutralBorder, line width=1pt,
        fill=White, minimum width=\boxW cm, minimum height=\boxH cm,
        align=center, text=TextMain, font=\bfseries\small, #1]
        (tmp) at (#2,#3) {};
  \node[align=center, text=TextMain, font=\bfseries\footnotesize]
        at ($(tmp.north)+(0,-0.3)$) {#4};
  \node[align=center, text=TextSub, font=\scriptsize\itshape]
        at ($(tmp.south)+(0,0.3)$) {#5};
}

\begin{tikzpicture}[
    scale=0.92,
    >={Stealth[length=6pt,width=5pt]},
    column/.style={
        rounded corners=5pt, draw=NeutralMid, line width=1pt,
        fill=NeutralColFill, minimum width=\colW cm, minimum height=\colH cm,
        inner sep=0pt, anchor=south west
    },
    perceptionColumn/.style={
        rounded corners=5pt, draw=AccentBlue, line width=1.6pt,
        fill=AccentBlueFill, minimum width=\colW cm, minimum height=\colH cm,
        inner sep=0pt, anchor=south west
    },
    titleBar/.style={
        rounded corners=4pt, draw=AccentBlue, line width=0.9pt,
        fill=AccentBlue, minimum width=\colW cm, minimum height=\titleH cm,
        align=center, text=White, font=\bfseries\small, anchor=south west
    },
    titleBarNeutral/.style={
        rounded corners=4pt, draw=NeutralMid, line width=0.9pt,
        fill=NeutralMid, minimum width=\colW cm, minimum height=\titleH cm,
        align=center, text=White, font=\bfseries\small, anchor=south west
    }
  ]
  % --- Column coordinates (5 Columns) ---
  \def\xA{0}
  \def\xB{\colW+\colGap}
  \def\xC{2*(\colW+\colGap)}
  \def\xD{3*(\colW+\colGap)}
  \def\xE{4*(\colW+\colGap)}
  \def\yTop{0}
  \def\yBot{-\colH}
  % --- Background Columns ---
  \node[column]            (C1) at (\xA, \yBot) {};
  \node[perceptionColumn]  (C2) at (\xB, \yBot) {};
  \node[column]            (C3) at (\xC, \yBot) {};
  \node[column]            (C4) at (\xD, \yBot) {};
  \node[column]            (C5) at (\xE, \yBot) {};
  % --- Column Header Titles ---
  \node[titleBarNeutral]   (T1) at (\xA, \yTop-\titleH) {1. SENSING \& INPUTS};
  \node[titleBar]          (T2) at (\xB, \yTop-\titleH) {2. PERCEPTION};
  \node[titleBarNeutral]   (T3) at (\xC, \yTop-\titleH) {3. KINEMATICS \& ALIGN};
  \node[titleBarNeutral]   (T4) at (\xD, \yTop-\titleH) {4. EVALUATION};
  \node[titleBarNeutral]   (T5) at (\xE, \yTop-\titleH) {5. COACHING \& HUD};
  % --- Modality dividers in Perception Column (Dashed hairlines) ---
  \draw[NeutralMid, line width=0.6pt, dash pattern=on 3pt off 3pt]
       ($(C2.north west)+(0.15,-2.5)$) -- ($(C2.north east)+(-0.15,-2.5)$);
  \draw[NeutralMid, line width=0.6pt, dash pattern=on 3pt off 3pt]
       ($(C2.north west)+(0.15,-4.9)$) -- ($(C2.north east)+(-0.15,-4.9)$);
  % Modality Label tags inside Perception
  \node[font=\scriptsize\bfseries\itshape, text=AccentBlue, anchor=east]
       at ($(C2.north east)+(-0.12,-0.95)$) {Body};
  \node[font=\scriptsize\bfseries\itshape, text=AccentBlue, anchor=east]
       at ($(C2.north east)+(-0.12,-3.0)$) {Face};
  \node[font=\scriptsize\bfseries\itshape, text=AccentBlue, anchor=east]
       at ($(C2.north east)+(-0.12,-5.3)$) {Audio};
  % ===========================================================================
  %  STAGE COMPONENT BOXES
  % ===========================================================================
  % --- Column 1 -- SENSING & INPUTS ---
  \compbox{\xA+\colW/2}{-1.75}{User Video / Webcam}{30 FPS RGB Frame Stream}
  \compbox{\xA+\colW/2}{-3.85}{Reference Template}{Exercise Angular Dataset}
  \compbox{\xA+\colW/2}{-5.95}{Safe-Max ROM Calibration}{2s Neutral + 5s Max ($P_{95}$)}
  % --- Column 2 -- PERCEPTION ---
  \compbox{\xB+\colW/2}{-1.75}{RTMW3D-L Model}{133 3D Keypoint Extraction}
  \compbox{\xB+\colW/2}{-3.85}{OpenFace 3.0 Engine}{Action Units (8 AUs) + Gaze}
  \compbox{\xB+\colW/2}{-5.95}{Whisper ASR Model}{Speech-to-Text Fallback}
  % --- Column 3 -- KINEMATICS & ALIGN ---
  \compbox{\xC+\colW/2}{-1.75}{Quaternion Kinematics}{Melax Robust included angle}
  \compbox{\xC+\colW/2}{-3.85}{Butterworth Filter}{4th-Order, 6.0 Hz Cutoff}
  \compbox{\xC+\colW/2}{-5.95}{Constrained DTW}{Sakoe-Chiba Band (15\% width)}
  % --- Column 4 -- EVALUATION ---
  \compbox{\xD+\colW/2}{-1.75}{6-Dimension Scorer}{ROM (25\%) + Stability (15\%)}
  \compbox{\xD+\colW/2}{-3.85}{Compensation Detector}{Shoulder Hiking + Trunk Lean}
  \compbox{\xD+\colW/2}{-5.95}{Fatigue Analyzer}{Jerk Ratio + ROM Degradation}
  % --- Column 5 -- COACHING & HUD ---
  \compbox{\xE+\colW/2}{-1.75}{LLM Coach \& RAG}{Clinical KB + Safety Guardrails}
  \compbox{\xE+\colW/2}{-3.85}{Edge-TTS Speech Synth}{vi-VN-HoaiMyNeural Voice}
  \compbox{\xE+\colW/2}{-5.95}{Visual HUD \& Dashboard}{2D skeleton overlay + ROM arcs}
  % ===========================================================================
  %  ROUTING / CONNECTION ARROWS
  % ===========================================================================

  % Main Flow Arrows (y = -3.2)
  \draw[NeutralMid, line width=1.5pt, -{Stealth[length=6pt,width=5pt]}]
       ($(C1.east)+(0,-3.2)$) -- ($(C2.west)+(0,-3.2)$);
  \draw[NeutralMid, line width=1.5pt, -{Stealth[length=6pt,width=5pt]}]
       ($(C2.east)+(0,-3.2)$) -- ($(C3.west)+(0,-3.2)$);
  \draw[NeutralMid, line width=1.5pt, -{Stealth[length=6pt,width=5pt]}]
       ($(C3.east)+(0,-3.2)$) -- ($(C4.west)+(0,-3.2)$);
  \draw[NeutralMid, line width=1.5pt, -{Stealth[length=6pt,width=5pt]}]
       ($(C4.east)+(0,-3.2)$) -- ($(C5.west)+(0,-3.2)$);
  % Arrow Labels
  \node[font=\scriptsize\bfseries, text=TextMain, fill=White, inner sep=1.5pt]
       at ($(C1.east)!0.5!(C2.west) + (0,0.25)$) {RGB feed};
  \node[font=\scriptsize\bfseries, text=TextMain, fill=White, inner sep=1.5pt]
       at ($(C2.east)!0.5!(C3.west) + (0,0.25)$) {Keypoints};
  \node[font=\scriptsize\bfseries, text=TextMain, fill=White, inner sep=1.5pt]
       at ($(C3.east)!0.5!(C4.west) + (0,0.25)$) {Aligned Angles};
  \node[font=\scriptsize\bfseries, text=TextMain, fill=White, inner sep=1.5pt]
       at ($(C4.east)!0.5!(C5.west) + (0,0.25)$) {6-D Scores};
  % Internal Module Sub-routing Arrows
  \draw[NeutralMid, line width=0.8pt, ->] ($(C1.east)+(0,-5.95)$) -- +(0.15,0) |- ($(C3.west)+(0,-5.95)$);
  \draw[NeutralMid, line width=0.8pt, ->] ($(C2.east)+(0,-1.75)$) -- ($(C3.west)+(0,-1.75)$);
  \draw[NeutralMid, line width=0.8pt, ->] ($(C2.east)+(0,-3.85)$) -- +(0.15,0) |- ($(C3.west)+(0,-4.5)$);
  \draw[NeutralMid, line width=0.8pt, ->] ($(C3.east)+(0,-1.75)$) -- ($(C4.west)+(0,-1.75)$);
  \draw[NeutralMid, line width=0.8pt, ->] ($(C3.east)+(0,-5.95)$) -- +(0.15,0) |- ($(C4.west)+(0,-5.95)$);
  \draw[NeutralMid, line width=0.8pt, ->] ($(C4.east)+(0,-1.75)$) -- ($(C5.west)+(0,-1.75)$);
  \draw[NeutralMid, line width=0.8pt, ->] ($(C4.east)+(0,-5.95)$) -- +(0.15,0) |- ($(C5.west)+(0,-5.95)$);
  % Active feedback loops (Pain/Fatigue alerts to Coach, red dashed line)
  \coordinate (FBstart) at ($(C5.south)+(0,-0.12)$);
  \coordinate (FBend)   at ($(C4.south)+(0,-0.12)$);
  \draw[AccentRed, line width=1.4pt, dash pattern=on 5pt off 4pt, -{Stealth[length=6pt,width=5pt]}]
        (FBstart) .. controls +(0,-1.3) and +(0,-1.3) .. (FBend);
  \node[font=\scriptsize\bfseries\itshape, text=AccentRed, fill=White, inner sep=2pt]
        at ($(FBstart)!0.5!(FBend)+(0,-1.1)$) {Pain \& Fatigue Loop $\rightarrow$ Dynamic Coach Adjust};
\end{tikzpicture}
 into IEEEtran figure*
%  Total visual width ≈ 18.1 cm with scale=0.92 (fits IEEE 2-col textwidth)
% =============================================================================

% --- Dimension constants (cm) ------------------------------------------------
\def\colW{3.5}
\def\colH{7.8}
\def\colGap{0.55}
\def\boxW{3.1}
\def\boxH{1.0}
\def\titleH{0.6}

% --- Component Box Helper Macro ----------------------------------------------
\newcommand{\compbox}[5][]{%
  \node[rounded corners=4pt, draw=NeutralBorder, line width=1pt,
        fill=White, minimum width=\boxW cm, minimum height=\boxH cm,
        align=center, text=TextMain, font=\bfseries\small, #1]
        (tmp) at (#2,#3) {};
  \node[align=center, text=TextMain, font=\bfseries\footnotesize]
        at ($(tmp.north)+(0,-0.3)$) {#4};
  \node[align=center, text=TextSub, font=\scriptsize\itshape]
        at ($(tmp.south)+(0,0.3)$) {#5};
}

\begin{tikzpicture}[
    scale=0.92,
    >={Stealth[length=6pt,width=5pt]},
    column/.style={
        rounded corners=5pt, draw=NeutralMid, line width=1pt,
        fill=NeutralColFill, minimum width=\colW cm, minimum height=\colH cm,
        inner sep=0pt, anchor=south west
    },
    perceptionColumn/.style={
        rounded corners=5pt, draw=AccentBlue, line width=1.6pt,
        fill=AccentBlueFill, minimum width=\colW cm, minimum height=\colH cm,
        inner sep=0pt, anchor=south west
    },
    titleBar/.style={
        rounded corners=4pt, draw=AccentBlue, line width=0.9pt,
        fill=AccentBlue, minimum width=\colW cm, minimum height=\titleH cm,
        align=center, text=White, font=\bfseries\small, anchor=south west
    },
    titleBarNeutral/.style={
        rounded corners=4pt, draw=NeutralMid, line width=0.9pt,
        fill=NeutralMid, minimum width=\colW cm, minimum height=\titleH cm,
        align=center, text=White, font=\bfseries\small, anchor=south west
    }
  ]
  % --- Column coordinates (5 Columns) ---
  \def\xA{0}
  \def\xB{\colW+\colGap}
  \def\xC{2*(\colW+\colGap)}
  \def\xD{3*(\colW+\colGap)}
  \def\xE{4*(\colW+\colGap)}
  \def\yTop{0}
  \def\yBot{-\colH}
  % --- Background Columns ---
  \node[column]            (C1) at (\xA, \yBot) {};
  \node[perceptionColumn]  (C2) at (\xB, \yBot) {};
  \node[column]            (C3) at (\xC, \yBot) {};
  \node[column]            (C4) at (\xD, \yBot) {};
  \node[column]            (C5) at (\xE, \yBot) {};
  % --- Column Header Titles ---
  \node[titleBarNeutral]   (T1) at (\xA, \yTop-\titleH) {1. SENSING \& INPUTS};
  \node[titleBar]          (T2) at (\xB, \yTop-\titleH) {2. PERCEPTION};
  \node[titleBarNeutral]   (T3) at (\xC, \yTop-\titleH) {3. KINEMATICS \& ALIGN};
  \node[titleBarNeutral]   (T4) at (\xD, \yTop-\titleH) {4. EVALUATION};
  \node[titleBarNeutral]   (T5) at (\xE, \yTop-\titleH) {5. COACHING \& HUD};
  % --- Modality dividers in Perception Column (Dashed hairlines) ---
  \draw[NeutralMid, line width=0.6pt, dash pattern=on 3pt off 3pt]
       ($(C2.north west)+(0.15,-2.5)$) -- ($(C2.north east)+(-0.15,-2.5)$);
  \draw[NeutralMid, line width=0.6pt, dash pattern=on 3pt off 3pt]
       ($(C2.north west)+(0.15,-4.9)$) -- ($(C2.north east)+(-0.15,-4.9)$);
  % Modality Label tags inside Perception
  \node[font=\scriptsize\bfseries\itshape, text=AccentBlue, anchor=east]
       at ($(C2.north east)+(-0.12,-0.95)$) {Body};
  \node[font=\scriptsize\bfseries\itshape, text=AccentBlue, anchor=east]
       at ($(C2.north east)+(-0.12,-3.0)$) {Face};
  \node[font=\scriptsize\bfseries\itshape, text=AccentBlue, anchor=east]
       at ($(C2.north east)+(-0.12,-5.3)$) {Audio};
  % ===========================================================================
  %  STAGE COMPONENT BOXES
  % ===========================================================================
  % --- Column 1 -- SENSING & INPUTS ---
  \compbox{\xA+\colW/2}{-1.75}{User Video / Webcam}{30 FPS RGB Frame Stream}
  \compbox{\xA+\colW/2}{-3.85}{Reference Template}{Exercise Angular Dataset}
  \compbox{\xA+\colW/2}{-5.95}{Safe-Max ROM Calibration}{2s Neutral + 5s Max ($P_{95}$)}
  % --- Column 2 -- PERCEPTION ---
  \compbox{\xB+\colW/2}{-1.75}{RTMW3D-L Model}{133 3D Keypoint Extraction}
  \compbox{\xB+\colW/2}{-3.85}{OpenFace 3.0 Engine}{Action Units (8 AUs) + Gaze}
  \compbox{\xB+\colW/2}{-5.95}{Whisper ASR Model}{Speech-to-Text Fallback}
  % --- Column 3 -- KINEMATICS & ALIGN ---
  \compbox{\xC+\colW/2}{-1.75}{Quaternion Kinematics}{Melax Robust included angle}
  \compbox{\xC+\colW/2}{-3.85}{Butterworth Filter}{4th-Order, 6.0 Hz Cutoff}
  \compbox{\xC+\colW/2}{-5.95}{Constrained DTW}{Sakoe-Chiba Band (15\% width)}
  % --- Column 4 -- EVALUATION ---
  \compbox{\xD+\colW/2}{-1.75}{6-Dimension Scorer}{ROM (25\%) + Stability (15\%)}
  \compbox{\xD+\colW/2}{-3.85}{Compensation Detector}{Shoulder Hiking + Trunk Lean}
  \compbox{\xD+\colW/2}{-5.95}{Fatigue Analyzer}{Jerk Ratio + ROM Degradation}
  % --- Column 5 -- COACHING & HUD ---
  \compbox{\xE+\colW/2}{-1.75}{LLM Coach \& RAG}{Clinical KB + Safety Guardrails}
  \compbox{\xE+\colW/2}{-3.85}{Edge-TTS Speech Synth}{vi-VN-HoaiMyNeural Voice}
  \compbox{\xE+\colW/2}{-5.95}{Visual HUD \& Dashboard}{2D skeleton overlay + ROM arcs}
  % ===========================================================================
  %  ROUTING / CONNECTION ARROWS
  % ===========================================================================

  % Main Flow Arrows (y = -3.2)
  \draw[NeutralMid, line width=1.5pt, -{Stealth[length=6pt,width=5pt]}]
       ($(C1.east)+(0,-3.2)$) -- ($(C2.west)+(0,-3.2)$);
  \draw[NeutralMid, line width=1.5pt, -{Stealth[length=6pt,width=5pt]}]
       ($(C2.east)+(0,-3.2)$) -- ($(C3.west)+(0,-3.2)$);
  \draw[NeutralMid, line width=1.5pt, -{Stealth[length=6pt,width=5pt]}]
       ($(C3.east)+(0,-3.2)$) -- ($(C4.west)+(0,-3.2)$);
  \draw[NeutralMid, line width=1.5pt, -{Stealth[length=6pt,width=5pt]}]
       ($(C4.east)+(0,-3.2)$) -- ($(C5.west)+(0,-3.2)$);
  % Arrow Labels
  \node[font=\scriptsize\bfseries, text=TextMain, fill=White, inner sep=1.5pt]
       at ($(C1.east)!0.5!(C2.west) + (0,0.25)$) {RGB feed};
  \node[font=\scriptsize\bfseries, text=TextMain, fill=White, inner sep=1.5pt]
       at ($(C2.east)!0.5!(C3.west) + (0,0.25)$) {Keypoints};
  \node[font=\scriptsize\bfseries, text=TextMain, fill=White, inner sep=1.5pt]
       at ($(C3.east)!0.5!(C4.west) + (0,0.25)$) {Aligned Angles};
  \node[font=\scriptsize\bfseries, text=TextMain, fill=White, inner sep=1.5pt]
       at ($(C4.east)!0.5!(C5.west) + (0,0.25)$) {6-D Scores};
  % Internal Module Sub-routing Arrows
  \draw[NeutralMid, line width=0.8pt, ->] ($(C1.east)+(0,-5.95)$) -- +(0.15,0) |- ($(C3.west)+(0,-5.95)$);
  \draw[NeutralMid, line width=0.8pt, ->] ($(C2.east)+(0,-1.75)$) -- ($(C3.west)+(0,-1.75)$);
  \draw[NeutralMid, line width=0.8pt, ->] ($(C2.east)+(0,-3.85)$) -- +(0.15,0) |- ($(C3.west)+(0,-4.5)$);
  \draw[NeutralMid, line width=0.8pt, ->] ($(C3.east)+(0,-1.75)$) -- ($(C4.west)+(0,-1.75)$);
  \draw[NeutralMid, line width=0.8pt, ->] ($(C3.east)+(0,-5.95)$) -- +(0.15,0) |- ($(C4.west)+(0,-5.95)$);
  \draw[NeutralMid, line width=0.8pt, ->] ($(C4.east)+(0,-1.75)$) -- ($(C5.west)+(0,-1.75)$);
  \draw[NeutralMid, line width=0.8pt, ->] ($(C4.east)+(0,-5.95)$) -- +(0.15,0) |- ($(C5.west)+(0,-5.95)$);
  % Active feedback loops (Pain/Fatigue alerts to Coach, red dashed line)
  \coordinate (FBstart) at ($(C5.south)+(0,-0.12)$);
  \coordinate (FBend)   at ($(C4.south)+(0,-0.12)$);
  \draw[AccentRed, line width=1.4pt, dash pattern=on 5pt off 4pt, -{Stealth[length=6pt,width=5pt]}]
        (FBstart) .. controls +(0,-1.3) and +(0,-1.3) .. (FBend);
  \node[font=\scriptsize\bfseries\itshape, text=AccentRed, fill=White, inner sep=2pt]
        at ($(FBstart)!0.5!(FBend)+(0,-1.1)$) {Pain \& Fatigue Loop $\rightarrow$ Dynamic Coach Adjust};
\end{tikzpicture}
 into IEEEtran figure*
%  Total visual width ≈ 18.1 cm with scale=0.92 (fits IEEE 2-col textwidth)
% =============================================================================

% --- Dimension constants (cm) ------------------------------------------------
\def\colW{3.5}
\def\colH{7.8}
\def\colGap{0.55}
\def\boxW{3.1}
\def\boxH{1.0}
\def\titleH{0.6}

% --- Component Box Helper Macro ----------------------------------------------
\newcommand{\compbox}[5][]{%
  \node[rounded corners=4pt, draw=NeutralBorder, line width=1pt,
        fill=White, minimum width=\boxW cm, minimum height=\boxH cm,
        align=center, text=TextMain, font=\bfseries\small, #1]
        (tmp) at (#2,#3) {};
  \node[align=center, text=TextMain, font=\bfseries\footnotesize]
        at ($(tmp.north)+(0,-0.3)$) {#4};
  \node[align=center, text=TextSub, font=\scriptsize\itshape]
        at ($(tmp.south)+(0,0.3)$) {#5};
}

\begin{tikzpicture}[
    scale=0.92,
    >={Stealth[length=6pt,width=5pt]},
    column/.style={
        rounded corners=5pt, draw=NeutralMid, line width=1pt,
        fill=NeutralColFill, minimum width=\colW cm, minimum height=\colH cm,
        inner sep=0pt, anchor=south west
    },
    perceptionColumn/.style={
        rounded corners=5pt, draw=AccentBlue, line width=1.6pt,
        fill=AccentBlueFill, minimum width=\colW cm, minimum height=\colH cm,
        inner sep=0pt, anchor=south west
    },
    titleBar/.style={
        rounded corners=4pt, draw=AccentBlue, line width=0.9pt,
        fill=AccentBlue, minimum width=\colW cm, minimum height=\titleH cm,
        align=center, text=White, font=\bfseries\small, anchor=south west
    },
    titleBarNeutral/.style={
        rounded corners=4pt, draw=NeutralMid, line width=0.9pt,
        fill=NeutralMid, minimum width=\colW cm, minimum height=\titleH cm,
        align=center, text=White, font=\bfseries\small, anchor=south west
    }
  ]
  % --- Column coordinates (5 Columns) ---
  \def\xA{0}
  \def\xB{\colW+\colGap}
  \def\xC{2*(\colW+\colGap)}
  \def\xD{3*(\colW+\colGap)}
  \def\xE{4*(\colW+\colGap)}
  \def\yTop{0}
  \def\yBot{-\colH}
  % --- Background Columns ---
  \node[column]            (C1) at (\xA, \yBot) {};
  \node[perceptionColumn]  (C2) at (\xB, \yBot) {};
  \node[column]            (C3) at (\xC, \yBot) {};
  \node[column]            (C4) at (\xD, \yBot) {};
  \node[column]            (C5) at (\xE, \yBot) {};
  % --- Column Header Titles ---
  \node[titleBarNeutral]   (T1) at (\xA, \yTop-\titleH) {1. SENSING \& INPUTS};
  \node[titleBar]          (T2) at (\xB, \yTop-\titleH) {2. PERCEPTION};
  \node[titleBarNeutral]   (T3) at (\xC, \yTop-\titleH) {3. KINEMATICS \& ALIGN};
  \node[titleBarNeutral]   (T4) at (\xD, \yTop-\titleH) {4. EVALUATION};
  \node[titleBarNeutral]   (T5) at (\xE, \yTop-\titleH) {5. COACHING \& HUD};
  % --- Modality dividers in Perception Column (Dashed hairlines) ---
  \draw[NeutralMid, line width=0.6pt, dash pattern=on 3pt off 3pt]
       ($(C2.north west)+(0.15,-2.5)$) -- ($(C2.north east)+(-0.15,-2.5)$);
  \draw[NeutralMid, line width=0.6pt, dash pattern=on 3pt off 3pt]
       ($(C2.north west)+(0.15,-4.9)$) -- ($(C2.north east)+(-0.15,-4.9)$);
  % Modality Label tags inside Perception
  \node[font=\scriptsize\bfseries\itshape, text=AccentBlue, anchor=east]
       at ($(C2.north east)+(-0.12,-0.95)$) {Body};
  \node[font=\scriptsize\bfseries\itshape, text=AccentBlue, anchor=east]
       at ($(C2.north east)+(-0.12,-3.0)$) {Face};
  \node[font=\scriptsize\bfseries\itshape, text=AccentBlue, anchor=east]
       at ($(C2.north east)+(-0.12,-5.3)$) {Audio};
  % ===========================================================================
  %  STAGE COMPONENT BOXES
  % ===========================================================================
  % --- Column 1 -- SENSING & INPUTS ---
  \compbox{\xA+\colW/2}{-1.75}{User Video / Webcam}{30 FPS RGB Frame Stream}
  \compbox{\xA+\colW/2}{-3.85}{Reference Template}{Exercise Angular Dataset}
  \compbox{\xA+\colW/2}{-5.95}{Safe-Max ROM Calibration}{2s Neutral + 5s Max ($P_{95}$)}
  % --- Column 2 -- PERCEPTION ---
  \compbox{\xB+\colW/2}{-1.75}{RTMW3D-L Model}{133 3D Keypoint Extraction}
  \compbox{\xB+\colW/2}{-3.85}{OpenFace 3.0 Engine}{Action Units (8 AUs) + Gaze}
  \compbox{\xB+\colW/2}{-5.95}{Whisper ASR Model}{Speech-to-Text Fallback}
  % --- Column 3 -- KINEMATICS & ALIGN ---
  \compbox{\xC+\colW/2}{-1.75}{Quaternion Kinematics}{Melax Robust included angle}
  \compbox{\xC+\colW/2}{-3.85}{Butterworth Filter}{4th-Order, 6.0 Hz Cutoff}
  \compbox{\xC+\colW/2}{-5.95}{Constrained DTW}{Sakoe-Chiba Band (15\% width)}
  % --- Column 4 -- EVALUATION ---
  \compbox{\xD+\colW/2}{-1.75}{6-Dimension Scorer}{ROM (25\%) + Stability (15\%)}
  \compbox{\xD+\colW/2}{-3.85}{Compensation Detector}{Shoulder Hiking + Trunk Lean}
  \compbox{\xD+\colW/2}{-5.95}{Fatigue Analyzer}{Jerk Ratio + ROM Degradation}
  % --- Column 5 -- COACHING & HUD ---
  \compbox{\xE+\colW/2}{-1.75}{LLM Coach \& RAG}{Clinical KB + Safety Guardrails}
  \compbox{\xE+\colW/2}{-3.85}{Edge-TTS Speech Synth}{vi-VN-HoaiMyNeural Voice}
  \compbox{\xE+\colW/2}{-5.95}{Visual HUD \& Dashboard}{2D skeleton overlay + ROM arcs}
  % ===========================================================================
  %  ROUTING / CONNECTION ARROWS
  % ===========================================================================

  % Main Flow Arrows (y = -3.2)
  \draw[NeutralMid, line width=1.5pt, -{Stealth[length=6pt,width=5pt]}]
       ($(C1.east)+(0,-3.2)$) -- ($(C2.west)+(0,-3.2)$);
  \draw[NeutralMid, line width=1.5pt, -{Stealth[length=6pt,width=5pt]}]
       ($(C2.east)+(0,-3.2)$) -- ($(C3.west)+(0,-3.2)$);
  \draw[NeutralMid, line width=1.5pt, -{Stealth[length=6pt,width=5pt]}]
       ($(C3.east)+(0,-3.2)$) -- ($(C4.west)+(0,-3.2)$);
  \draw[NeutralMid, line width=1.5pt, -{Stealth[length=6pt,width=5pt]}]
       ($(C4.east)+(0,-3.2)$) -- ($(C5.west)+(0,-3.2)$);
  % Arrow Labels
  \node[font=\scriptsize\bfseries, text=TextMain, fill=White, inner sep=1.5pt]
       at ($(C1.east)!0.5!(C2.west) + (0,0.25)$) {RGB feed};
  \node[font=\scriptsize\bfseries, text=TextMain, fill=White, inner sep=1.5pt]
       at ($(C2.east)!0.5!(C3.west) + (0,0.25)$) {Keypoints};
  \node[font=\scriptsize\bfseries, text=TextMain, fill=White, inner sep=1.5pt]
       at ($(C3.east)!0.5!(C4.west) + (0,0.25)$) {Aligned Angles};
  \node[font=\scriptsize\bfseries, text=TextMain, fill=White, inner sep=1.5pt]
       at ($(C4.east)!0.5!(C5.west) + (0,0.25)$) {6-D Scores};
  % Internal Module Sub-routing Arrows
  \draw[NeutralMid, line width=0.8pt, ->] ($(C1.east)+(0,-5.95)$) -- +(0.15,0) |- ($(C3.west)+(0,-5.95)$);
  \draw[NeutralMid, line width=0.8pt, ->] ($(C2.east)+(0,-1.75)$) -- ($(C3.west)+(0,-1.75)$);
  \draw[NeutralMid, line width=0.8pt, ->] ($(C2.east)+(0,-3.85)$) -- +(0.15,0) |- ($(C3.west)+(0,-4.5)$);
  \draw[NeutralMid, line width=0.8pt, ->] ($(C3.east)+(0,-1.75)$) -- ($(C4.west)+(0,-1.75)$);
  \draw[NeutralMid, line width=0.8pt, ->] ($(C3.east)+(0,-5.95)$) -- +(0.15,0) |- ($(C4.west)+(0,-5.95)$);
  \draw[NeutralMid, line width=0.8pt, ->] ($(C4.east)+(0,-1.75)$) -- ($(C5.west)+(0,-1.75)$);
  \draw[NeutralMid, line width=0.8pt, ->] ($(C4.east)+(0,-5.95)$) -- +(0.15,0) |- ($(C5.west)+(0,-5.95)$);
  % Active feedback loops (Pain/Fatigue alerts to Coach, red dashed line)
  \coordinate (FBstart) at ($(C5.south)+(0,-0.12)$);
  \coordinate (FBend)   at ($(C4.south)+(0,-0.12)$);
  \draw[AccentRed, line width=1.4pt, dash pattern=on 5pt off 4pt, -{Stealth[length=6pt,width=5pt]}]
        (FBstart) .. controls +(0,-1.3) and +(0,-1.3) .. (FBend);
  \node[font=\scriptsize\bfseries\itshape, text=AccentRed, fill=White, inner sep=2pt]
        at ($(FBstart)!0.5!(FBend)+(0,-1.1)$) {Pain \& Fatigue Loop $\rightarrow$ Dynamic Coach Adjust};
\end{tikzpicture}

    \caption{\system{} system architecture. Five pipeline stages process data from Sensing \& Inputs (camera, reference template, Safe-Max ROM calibration) through Perception (RTMW3D-L whole-body 3D pose, OpenFace~3.0 AUs, Whisper ASR) and Kinematics \& Alignment (quaternion angles, Butterworth filtering, constrained DTW) to Evaluation (6-D scoring, compensation detection, fatigue analysis) and Coaching \& HUD (LLM coach with RAG safety guardrails, Edge-TTS, visual overlay). Red dashed feedback arcs enable real-time pain/fatigue alerts to dynamically adjust coaching.}
    \label{fig:architecture}
\end{figure*}

The architecture supports graceful degradation: if the LLM API is unreachable, the coaching module uses rule-based text feedback; the face analysis pipeline automatically falls back from OpenFace~3.0 to geometric analysis based on available dependencies.

\subsection{Whole-Body 3D Pose Estimation}
\label{subsec:pose3d}

We adopt RTMW3D-L~\cite{ref_rtmw3d} from MMPose~\cite{ref_mmpose} for real-time whole-body 3D pose estimation. RTMW3D-L employs a top-down Cascade Top-Down framework with a CSPNeXt backbone, CSPNeXtPAFPN neck, and RTMWHead that predicts both 2D heatmaps and 3D relative coordinates.

The model produces 133 keypoints directly from RGB images:
\begin{equation}
    \mathbf{P}_{3D} = f_{\text{RTMW3D-L}}(\mathbf{I}) \in \mathbb{R}^{133 \times 3}
    \label{eq:rtmw3d}
\end{equation}
where the output follows COCO-WholeBody format: body joints (indices 0--16, 17 points), foot keypoints (17--22, 6 points), face landmarks (23--90, 68 points in dlib convention), left hand (91--111, 21 points), and right hand (112--132, 21 points).

The model is pre-trained on the COCO-WholeBody dataset and fine-tuned on the Cocktail-14 multi-person 3D dataset. For inference we use the MMPose \texttt{inference\_topdown} API with PyTorch backend, achieving 129.7~FPS on an NVIDIA RTX 5880 Ada Generation GPU.

For joint angle computation we use the body keypoints (indices 0--16) in COCO format. Given three consecutive joints $\mathbf{a}$, $\mathbf{b}$, $\mathbf{c}$ forming an angle at $\mathbf{b}$, the dot-product angle is:
\begin{equation}
    \theta_{\text{dot}} = \arccos\left(\frac{(\mathbf{a} - \mathbf{b}) \cdot (\mathbf{c} - \mathbf{b})}{\|\mathbf{a} - \mathbf{b}\| \|\mathbf{c} - \mathbf{b}\|}\right) \cdot \frac{180}{\pi}
    \label{eq:angle_dot}
\end{equation}

The system computes angles for 8 bilateral joint pairs: shoulders, elbows, hips, and knees.

\subsection{Quaternion-Based Joint Angle Computation}
\label{subsec:quaternion}

Traditional dot-product angle computation is limited to the plane defined by three joints. We implement quaternion-based joint angles for multi-plane accuracy~\cite{ref_quaternion_angles}, using the Melax~(1998) robust formulation.

Given two limb segments with direction vectors $\mathbf{v}_1$ and $\mathbf{v}_2$, we compute:
\begin{equation}
    s = \sqrt{2(1 + \mathbf{v}_1 \cdot \mathbf{v}_2)}, \quad w = \frac{s}{2}, \quad \theta = 2 \arccos(w)
    \label{eq:quaternion}
\end{equation}

This sqrt-based formulation avoids numerical issues of direct $\arccos$ near 0 and 180 degrees. The system also implements ISB-convention clinical angles~\cite{ref_isb}: for flexion joints, $\theta_{\text{clinical}} = 180 - \theta_{\text{included}}$ (0 = full extension).

\begin{figure}[htbp]
    \centering
    % Figure 2: Safe-Max ROM Calibration Profile
\begin{tikzpicture}
\begin{axis}[
    width=8cm, height=5cm,
    xlabel={Time (seconds)},
    ylabel={Joint Angle (degrees)},
    xmin=0, xmax=5,
    ymin=20, ymax=160,
    grid=major,
    legend pos=south east,
    legend style={font=\footnotesize},
    tick label style={font=\footnotesize},
    label style={font=\small}
]

% Raw signal with jitter
\addplot[color=LightGray, thick, domain=0:5, samples=200]
    {80 + 50*sin(deg(2*x)) + 3*sin(deg(20*x)) + 2*cos(deg(35*x))};
\addlegendentry{Raw Signal}

% Filtered signal (smooth)
\addplot[color=RoyalBlue, very thick, domain=0:5, samples=100]
    {80 + 45*sin(deg(2*x))};
\addlegendentry{Filtered Signal}

% Upper outlier boundary (+2σ)
\addplot[color=gray, dashed, domain=0:5, samples=100]
    {80 + 45*sin(deg(2*x)) + 8};
\addlegendentry{Outlier $\pm 2\sigma$}

% Lower outlier boundary (-2σ)
\addplot[color=gray, dashed, domain=0:5, samples=100]
    {80 + 45*sin(deg(2*x)) - 8};

% P95 Target ROM Threshold
\addplot[color=BrickRed, dashed, ultra thick, domain=0:5] {122};
\addlegendentry{P95 Target}

% Annotation for P95
\node[anchor=west, font=\footnotesize] at (axis cs:3.2,122) {$\theta_{\text{user\_max}}$ (P95)};

\end{axis}
\end{tikzpicture}

    \caption{Safe-Max ROM calibration profile showing raw tracking signal with elderly tremor (light gray), median-filtered signal (royal blue), outlier rejection boundaries ($\pm 2\sigma$), and the P95 target threshold ($\theta_{\text{user\_max}}$) derived from the 95th percentile of filtered peak angles.}
    \label{fig:safemax_calibration}
\end{figure}

\subsection{SPARC Smoothness Metric}
\label{subsec:sparc}

Movement smoothness is a clinically validated marker of motor control quality~\cite{ref_sparc}. We adopt the Spectral Arc Length (SPARC) metric following Balasubramanian et al.~(2012).

SPARC computes the arc length of the normalized Fourier magnitude spectrum of the velocity profile up to a cutoff frequency $\omega_c$:
\begin{equation}
    \text{SPARC} = -\int_0^{\omega_c} \sqrt{\left(\frac{d\bar{\omega}}{\omega_c}\right)^2 + \left(\frac{d\hat{V}(\omega)}{d\omega}\right)^2} d\omega
    \label{eq:sparc}
\end{equation}
where $\hat{V}(\omega)$ is the peak-normalized magnitude spectrum and $\omega_c$ is determined adaptively (threshold = 0.05).

SPARC lies in $[-6, 0]$ with more negative values indicating smoother movements. We normalize to $[0, 100]$: $\text{score} = (\text{SPARC} + 6) / 6 \times 100$. The smoothness score combines SPARC (60\%) and Log-Dimensionless Jerk (LDLJ~\cite{ref_ldlj}, 40\%) for a robust composite.

\begin{figure}[htbp]
    \centering
    % Figure 4: SPARC Frequency-Domain Smoothness Comparison
\begin{tikzpicture}
\begin{axis}[
    width=8cm, height=5cm,
    xlabel={Frequency (Hz)},
    ylabel={Normalized Amplitude $\hat{M}(\omega)$},
    xmin=0, xmax=15,
    ymin=0, ymax=1.2,
    grid=major,
    legend pos=north east,
    legend style={font=\footnotesize},
    tick label style={font=\footnotesize},
    label style={font=\small}
]

% Curve A: Smooth / Healthy Movement (single sharp peak, rapid decay)
\addplot[color=ForestGreen, very thick, domain=0:15, samples=200]
    {exp(-5*x^2)};
\addlegendentry{Smooth (Healthy)}

% Curve B: Tremor / Impaired Movement (main peak + harmonic spikes)
\addplot[color=BrickRed, very thick, domain=0:15, samples=200]
    {exp(-3*x^2) + 0.15*exp(-20*(x-4)^2) + 0.12*exp(-20*(x-6.5)^2) + 0.08*exp(-20*(x-9)^2)};
\addlegendentry{Tremor (Impaired)}

% Threshold line (Y = 0.05)
\addplot[color=black, dashed, domain=0:15] {0.05};
\node[anchor=west, font=\footnotesize] at (axis cs:12,0.05) {Threshold};

% Cutoff frequency annotation (vertical line at ω_c ≈ 2.5 Hz for tremor curve)
\draw[gray, dashed] (axis cs:2.5,0) -- (axis cs:2.5,0.05);
\node[anchor=north, font=\footnotesize] at (axis cs:2.5,-0.05) {$\omega_c$};

\end{axis}
\end{tikzpicture}

    \caption{SPARC frequency-domain smoothness comparison. The smooth healthy movement (green) exhibits a single sharp peak with rapid decay, while the tremor/impaired movement (red) shows multiple high-frequency harmonic spikes. The cutoff frequency $\omega_c$ is set where the spectrum crosses the 0.05 threshold.}
    \label{fig:sparc_spectrum}
\end{figure}

\subsection{Multi-Indicator Facial State Detection}
\label{subsec:pain}

Our face analysis pipeline detects five facial states (Normal, Pain, Fatigue, Exhaustion, Boredom) in priority order: Pain, then Exhaustion, then Fatigue, then Boredom, then Normal.

\textbf{Face Detection}: Face keypoints come directly from RTMW3D's 133-keypoint output (COCO-WholeBody format, indices 23--90), providing 68 face landmarks in dlib convention. No separate face detection model is necessary---face, body, and hand keypoints are produced in a single inference pass, ensuring temporal consistency and reducing computational overhead.

\textbf{Action Unit Detection}: OpenFace~3.0 detects 8 Action Units (AU1, AU2, AU4, AU6, AU9, AU12, AU25, AU26) and 8 emotions (AffectNet classes). When OpenFace~3.0 is unavailable, the system falls back to geometric analysis of Face Mesh landmarks.

\textbf{Pain Intensity Estimation}: We adopt a PSPI-inspired index~\cite{ref_pspi} adapted to the available Action Units:
\begin{equation}
    \widetilde{\text{PSPI}} = \text{AU4} + 2 \cdot \text{AU6} + \text{AU9} + 2 \cdot \text{AU43}_{\text{approx}}
    \label{eq:pspi}
\end{equation}
where $\text{AU43}_{\text{approx}}$ (eye closure) is estimated from the Eye Aspect Ratio (EAR~\cite{ref_ear}). Intensity levels: None (0), Mild (1), Moderate (3), Strong (6), Severe (10).

\textbf{Fatigue Detection}: A composite score combines four weighted indicators:
\begin{equation}
    F = 0.35\,P_c + 0.25\,R_b + 0.25\,D_b + 0.15\,R_y
    \label{eq:fatigue}
\end{equation}
where $P_c$ is PERCLOS~\cite{ref_perclos} (\% eye closure over 60~s), $R_b$ is blink rate (normal: 15--20/min), $D_b$ is mean blink duration, and $R_y$ is yawn frequency. Thresholds: Alert ($F<0.3$), Mild ($F<0.5$), Fatigued ($F<0.7$), Exhaustion ($F\geq 0.85$).

\textbf{Engagement Detection}: Based on Whitehill et al.~(2014)~\cite{ref_engagement}:
\begin{equation}
    E = \alpha \cdot \text{AU12} - \gamma \cdot \text{AU1} + \delta
    \label{eq:engagement}
\end{equation}
with boredom index from D'Mello \& Graesser~(2010) based on low AU variability and prolonged neutral state.

\textbf{Blink and Yawn Detection}: Blinks are detected as brief EAR drops ($<$300ms) below threshold (0.2). Yawns are detected from AU25 (Lips Part) + AU26 (Jaw Drop) active simultaneously for $\geq$2 seconds.

\subsection{Voice-Interactive LLM Coaching}
\label{subsec:llm}

The coaching system integrates three components:
\begin{itemize}[itemsep=2pt, topsep=3pt, parsep=0pt]
    \item \textbf{LLM}: Large Language Models with RAG grounding from a clinical rehabilitation knowledge base covering exercise guidelines, contraindications, and region-specific healthcare context. The LLM receives structured context including exercise name, current phase, joint angle deviations, pain/fatigue levels, and user profile (age, ROM limits, medical conditions).
    \item \textbf{TTS}: Edge-TTS with neural voices for natural speech synthesis in the user's preferred language, converting LLM-generated coaching text to spoken feedback.
    \item \textbf{Safety}: Guardrails block harmful advice (e.g., ``ignore pain'', ``push through'') and check contraindications based on user profile.
\end{itemize}

A priority system ensures pain alerts are delivered immediately, followed by fatigue warnings, performance feedback, and encouragement. A 5-second cooldown between LLM calls prevents excessive API usage.

\subsection{Advanced Scoring Framework}
\label{subsec:scoring}

The scoring system evaluates exercises across six dimensions with clinically-inspired weights:

\begin{table}[htbp]
    \caption{Scoring Dimensions and Weights}
    \label{tab:scoring}
    \centering
    \begin{tabular}{lcl}
        \toprule
        \textbf{Dimension} & \textbf{Wt.} & \textbf{Metric} \\
        \midrule
        ROM Accuracy & 25\% & Max angle vs target \\
        Stability & 15\% & Angle variance \\
        Flow & 20\% & Velocity smoothness \\
        Symmetry & 15\% & L-R angle diff. \\
        Compensation & 15\% & Temporal detection \\
        Smoothness & 10\% & SPARC metric \\
        \bottomrule
    \end{tabular}
\end{table}

\textbf{Constrained DTW}: We apply the Sakoe-Chiba band constraint to prevent pathological alignments:
\begin{equation}
    w = \max\left(\lfloor \max(n,m) \times 0.1 \rfloor, |n - m|\right)
    \label{eq:dtw_constraint}
\end{equation}

\begin{figure}[htbp]
    \centering
    % Figure 3: Sakoe-Chiba Constrained DTW Warping Path
\begin{tikzpicture}[scale=0.4]

% Grid parameters
\def\N{12}
\def\M{12}
\def\w{2.5}

% Draw grid
\draw[step=1, gray, very thin] (0,0) grid (\N,\M);

% Draw Sakoe-Chiba constraint band (shaded diagonal corridor)
\fill[LightBlue, opacity=0.3]
    (0,0) -- (\w,0) -- (\N,\M-\w) -- (\N,\M) -- (\N-\w,\M) -- (0,\w) -- cycle;

% Draw optimal warping path (red line with steps)
\draw[BrickRed, line width=1.5pt]
    (0,0) -- (1,0) -- (1,1) -- (2,1) -- (2,2) -- (3,3) -- (4,3) --
    (5,4) -- (6,5) -- (7,6) -- (8,7) -- (9,8) -- (10,9) -- (11,10) --
    (11,11) -- (12,12);

% Axes
\draw[->, thick] (0,0) -- (\N+0.5,0) node[below right, xshift=2pt, font=\footnotesize] {Reference};
\draw[->, thick] (0,0) -- (0,\M+0.5) node[above left, yshift=2pt, font=\footnotesize] {User};

% Constraint width annotation (two-headed arrow)
\draw[<->, thick] (6,6) -- (6+\w,6+\w) node[midway, above right, font=\footnotesize] {$2w+1$};

% Corner markers
\fill (0,0) circle (3pt);
\fill (\N,\M) circle (3pt);

\end{tikzpicture}

    \caption{Sakoe-Chiba constrained DTW warping path. The shaded diagonal corridor (light blue) represents the constraint band of width $2w+1$, limiting the search space. The optimal warping path (red) aligns reference and user sequences monotonically while remaining within the constraint boundary.}
    \label{fig:dtw_warping}
\end{figure}

\textbf{Compensation Detection}: Temporal analysis of joint angle sequences identifies three types of compensatory movements: shoulder hiking (threshold: 5\% of frame height), trunk lean (threshold: 15$^\circ$), and hip shift (threshold: 6\% of frame height). Severity is computed as $\min(1, \text{max\_value} / (2 \cdot \text{threshold}))$.

\subsection{Cross-Skeleton Evaluation Framework}
\label{subsec:skeleton_mapping}

To enable evaluation against public benchmarks with different skeleton formats, we implement a comprehensive skeleton mapping system. The framework maps between three common skeleton types: MediaPipe (33 landmarks), SMPL-24 (24 joints), and Kinect~v2 (25 joints, used in UI-PRMD~\cite{ref_uiprmd}).

For cross-skeleton comparison, we define a common 14-joint evaluation subset (pelvis, spine, neck, head, and bilateral shoulder, elbow, wrist, hip, knee, ankle) that exists in all three skeletons. The mapping functions handle derived joints (e.g., pelvis = average of left/right hip) and index remapping.

The framework implements two standard MPJPE protocols so that future ground-truth evaluation can follow established conventions: Protocol~1 (root-aligned mean per-joint L2 distance) and Protocol~2 (PA-MPJPE, Procrustes-aligned via SVD over translation, rotation, and scale). It also provides an ICC(3,1) reliability routine (two-way mixed model, single measures) following Koo \& Li~(2016)~\cite{ref_icc}. In this paper these protocols are exercised only for self-consistency profiling; applying them against annotated ground truth is left to future work.
